\documentclass[11pt]{article}
\usepackage[T1]{fontenc}\usepackage[utf8]{inputenc}
\usepackage[margin=1.05in]{geometry}
\usepackage{amsmath,amssymb,amsthm,mathtools,lmodern,microtype}
\usepackage[authoryear,round]{natbib}\usepackage{xurl}
\usepackage[colorlinks=true,allcolors=blue]{hyperref}
\newcommand{\workbenchurl}{https://github.com/KokunoYumeto/schemes-and-moduli-workbench}
\newcommand{\recordurl}{https://github.com/KokunoYumeto/mathematics-workbenches}
\newcommand{\OO}{\mathcal O}\newcommand{\ZZ}{\mathbf Z}
\newcommand{\FF}{\mathbf F}\newcommand{\PP}{\mathbf P}
\newcommand{\et}{\mathrm{et}}\newcommand{\crys}{\mathrm{crys}}
\newcommand{\st}{\mathrm{st}}\newcommand{\tors}{\mathrm{tors}}
\DeclareMathOperator{\Spec}{Spec}\DeclareMathOperator{\Frac}{Frac}
\DeclareMathOperator{\Hom}{Hom}\DeclareMathOperator{\Pic}{Pic}
\newtheorem{theorem}{Theorem}[section]
\newtheorem{proposition}[theorem]{Proposition}\newtheorem{lemma}[theorem]{Lemma}
\newtheorem{corollary}[theorem]{Corollary}
\theoremstyle{definition}\newtheorem{definition}[theorem]{Definition}
\theoremstyle{remark}\newtheorem{remark}[theorem]{Remark}
\title{Grothendieck's lifting question for crystalline torsion,\\
after Bhatt--Morrow--Scholze}
\author{}\date{}
\begin{document}\maketitle
\begin{abstract}
Grothendieck proposed expressing a stable crystalline torsion condition through the $p$-adic cohomology of a smooth characteristic-zero lift. Ogus identifies the Frobenius-bijective part of integral crystalline cohomology with integral $p$-adic \'etale cohomology of the special fibre. Li--Liu identify the specialization kernel and give its sharp small-ramification bounds. These results yield a precise sufficient range for the desired vanishing. At the first boundary, we give explicit integral equations for the K3 cover and elliptic torsion section in the Bhatt--Morrow--Scholze construction. A smooth projective surface over $\ZZ_2$ then has torsion-free generic second \'etale cohomology and a one-dimensional stable part of special crystalline $2$-torsion. Its last hypersurface is chosen by finite-field Bertini. The counterexample phenomenon and the comparison theorems are known results; the purpose here is a complete account of the integral ingredients and their relation to Grothendieck's question. The mathematics and text were produced and checked by GPT-6 Astra Pro in ChatGPT Pro mode, including a separate statement-only reconstruction of both main assertions; no human has checked them.
\end{abstract}

\section{The lifting question and the two conclusions}

In his discussion of stable crystalline torsion, Grothendieck considers the case
\begin{quote}
``Lorsque X se rel\`eve (avec sa polarisation) en caract\'eristique nulle''
\end{quote}
and asks for an expression of the relevant condition in terms of the $p$-adic cohomology of the lift \citep[\S8.5.10, archive sheet~9, scan image~10]{Grothendieck}. The passage contains handwritten qualifications and cancellations. We address the precise proposed implication: does the absence of generic $p$-adic \'etale torsion imply the absence of stable crystalline $p$-torsion on the special fibre?

There is a positive answer in a strict ramification range, and a negative answer without that restriction. Both facts are consequences of established comparison theorems and constructions. It is essential to distinguish the special and generic \'etale groups: the specialization map goes from the former to the latter and can have a torsion kernel.

Let $k$ be perfect of characteristic $p>0$, put $W=W(k)$, and let $Z/k$ be smooth and projective. Set
\[
 C^j(Z)=H^j_{\crys}(Z/W),\qquad
 T^j(Z)=C^j(Z)_{\tors}.
\]
Frobenius is semilinear for the Witt-vector automorphism. For a finite-dimensional $k$-space with such an endomorphism, its stable part is the intersection of all Frobenius images; Frobenius is bijective on that part. In particular $T^j(Z)[p]_{\st}$ has this meaning. For finite $W$-modules the integral stable part is defined by compatible stable parts modulo powers of $p$.

Let $K$ be a complete discretely valued field of characteristic zero with residue field $k$, absolute ramification index $e$, and uniformizer $\pi$. Let $\mathcal X/\OO_K$ be smooth and proper with special fibre $Z$. After compatible choices of geometric points, write
\[
 s^j:H^j_{\et}(Z_{\bar k},\ZZ_p)
       \longrightarrow H^j_{\et}(\mathcal X_{\bar K},\ZZ_p)
\]
for specialization. Put $\mathfrak S=W(k)[[u]]$, with $u\mapsto\pi$ and $\varphi(u)=u^p$, and let
\[
 M^j=H^j_\Delta(\widehat{\mathcal X}/\mathfrak S)[u^\infty].
\]
Here $H^j_\Delta$ is Breuil--Kisin prismatic cohomology and $[u^\infty]$ denotes the submodule killed by some power of $u$. This module is of finite length. The quotient $M^j/uM^j$ inherits Frobenius.

\begin{theorem}[Integral obstruction and the valid range]\label{thm:obstruction}
For $j\geq1$ there are natural identifications
\begin{align}
 T^j(Z_{\bar k})[p]_{\st}
  &\simeq H^j_{\et}(Z_{\bar k},\ZZ_p)[p]
                          \otimes_{\FF_p}\bar k,\label{eq:stable}\\
 \ker s^j
  &\simeq\left((M^j/uM^j)\otimes_{W(k)}W(\bar k)\right)^{\varphi=1}.
                                                    \label{eq:kernel}
\end{align}
The first is Galois equivariant. The second is the specialization-kernel theorem of Li--Liu. Its consequences are:
\begin{enumerate}
\item If $e(j-1)<p-1$, then $s^j$ is injective.
\item If $e(j-1)<2(p-1)$, then $p^{j-1}\ker s^j=0$.
\item If $e(j-1)=p-1$, then $p\ker s^j=0$.
\end{enumerate}
Consequently, if $H^j_{\et}(\mathcal X_{\bar K},\ZZ_p)$ is torsion-free and $e(j-1)<p-1$, then $T^j(Z_{\bar k})[p]_{\st}=0$. For $\OO_K=W(k)$ this sufficient range is $j<p$.
\end{theorem}

Equation~\eqref{eq:stable} is Ogus's integral stable comparison, including torsion \citep[Theorem~4.9, Remark~4.10, Corollary~4.11, pp.~314--315]{Ogus}. Equation~\eqref{eq:kernel} and the three bounds are \citet[Theorem~4.14 and Corollary~4.15, author-version pp.~27--28]{LiLiu}. We give the reduction and hypotheses in the next section; no new prismatic comparison is claimed.

\begin{theorem}[An explicit model for the boundary counterexample]\label{thm:surface}
In $\PP^5_{\ZZ_2}$, with coordinates $(x_1,x_2,x_3,y_1,y_2,y_3)$, let $\mathcal Y$ be the intersection
\begin{equation}\label{eq:quadrics}
 x_i^2+x_jx_k+y_i^2-x_iy_i=0,
 \qquad \{i,j,k\}=\{1,2,3\},
\end{equation}
and let $\tau(x_i,y_i)=(x_i,x_i-y_i)$. Let
\[
 \mathcal E:\ y^2+xy=x^3+1,
 \qquad P=[2u:-u:8],
\]
where $u\in\ZZ_2$ is the unique odd root of $u^3+u^2+64=0$. The quotient
\[
 \mathcal D=(\mathcal Y\times\mathcal E)/
                     \langle(\tau,t_P)\rangle
\]
is smooth and projective of relative dimension three. It contains a smooth projective ample hypersurface $\mathcal X/\ZZ_2$ whose special fibre $X_0$ is geometrically connected and satisfies
\[
 H^2_{\et}(\mathcal X_{\overline{\mathbb Q}_2},\ZZ_2)
       \text{ is torsion-free},
 \qquad
 T^2(X_{0,\overline{\FF}_2})[2]_{\st}
       \simeq\overline{\FF}_2.
\]
Thus the unrestricted lifting implication fails already at $e=1$ and $j=p=2$, the first boundary of Theorem~\ref{thm:obstruction}.
\end{theorem}

\citet[Theorem~2.1 and Proposition~2.2, pp.~236--238]{BMS} already establish the counterexample phenomenon. Their argument uses an Enriques quotient of a K3 surface and a two-torsion section which disappears on the special fibre. Theorem~\ref{thm:surface} supplies explicit integral equations for those ingredients and proves the stable torsion assertion directly from fundamental groups and the coefficient exact sequence. The existence of the final smooth hypersurface is supplied by \citet[Theorem~1.1]{Poonen}; no numerical degree or unprovided equation is asserted.

The difficulty is integral: a rational comparison cannot detect the finite specialization kernel. The proof first identifies stable crystalline torsion with special-fibre \'etale torsion, where Li--Liu's bound applies. For the example, a double-cover character survives on the special fibre but is absorbed into an elliptic isogeny on the generic fibre. The coefficient exact sequence converts this difference in first cohomology into second-cohomology torsion. An ample hypersurface preserves the fundamental groups and produces the surface. Smooth projectivity supplies finiteness, Lefschetz and finite-field Bertini at precisely these steps.

\subsection*{Recent neighbouring questions}

Crystalline torsion also enters the obstruction theory of lifting varieties. \citet[Theorem~A and \S1.4]{Zou} states a criterion using closed forms, numerical Hodge--de Rham degeneration and torsion-freeness in the adjacent crystalline degree, together with symplectic examples separating existence of a lift from unobstructedness of its deformation functor. That preprint acknowledges computations with ChatGPT~6 Astra and~6.1 Sol. Its obstruction groups concern varying the lift; our calculation concerns the torsion comparison for a lift already constructed. Its results are contextual and are not inputs here.

\citet[Theorem~1.1 and \S1.2, pp.~1--3]{OpenAI} studies specialization of rational Hodge classes on abelian varieties and their pairings with divisors in \'etale and crystalline realizations. Such rational pairings address a different part of the good-reduction comparison: they cannot determine the finite kernel in \eqref{eq:kernel}. No theorem from that preprint is used in the present argument.

\section{Stable extraction and specialization}

Stable extraction isolates the part detected by $p$-adic characters; we now apply the two published comparisons to that part.

\begin{lemma}\label{lem:stable}
For finite $W(k)$-modules equipped with semilinear Frobenius, stable extraction is exact and commutes with taking the kernel of multiplication by $p$. Over an algebraically closed field, Ogus's comparison yields
\[
 C^j(Z)_{\st}\simeq H^j_{\et}(Z,\ZZ_p)\otimes_{\ZZ_p}W(k).
\]
\end{lemma}
\begin{proof}
First let the module be killed by $p^a$. It has finite length. Frobenius images and kernels stabilize, and the usual Fitting argument, applied to the Frobenius twist at each iterate, gives a direct sum of a bijective part and a nilpotent part. There are no Frobenius-compatible homomorphisms from a bijective part to a nilpotent part or conversely. This proves exactness of the stable summand on finite-length modules; it is also \citet[Lemma~4.8, p.~314]{Ogus}.

For a finite $W$-module, apply this construction modulo $p^a$ and pass to the inverse limit. The transition maps between stable summands are surjective by finite-length exactness, so their first derived inverse limit vanishes. For a submodule the induced $p$-adic filtration is equivalent to its own filtration by Artin--Rees. Thus an exact sequence of finite modules remains exact after stable extraction. Apply this to the kernel and image sequences of multiplication by $p$ to obtain
\[
 (C[p])_{\st}=C_{\st}[p].
\]
This argument also explains the inverse-limit passage in \citet[Remark~4.10]{Ogus}.

For the geometric assertion, a projective embedding of $Z$ into $\PP^N_k$ gives an embedding into the smooth $W_a(k)$-scheme $\PP^N_{W_a(k)}$ for every $a$. Thus the embedding hypothesis of Ogus's Theorem~4.9 is satisfied. That theorem identifies the stable crystalline part with ordinary infinitesimal cohomology; his Corollary~4.11 identifies its integral inverse limit with the displayed \'etale tensor product. The compatibility and exactness just proved retain torsion in that identification.
\end{proof}

\begin{proof}[Proof of Theorem~\ref{thm:obstruction}]
Apply Lemma~\ref{lem:stable} over $\bar k$. Since $W(\bar k)$ is flat over $\ZZ_p$, the kernel of $p$ on the \'etale tensor product is the tensor product of the kernel of $p$. The equality $C^j[p]=T^j[p]$ gives \eqref{eq:stable}. All maps are natural under automorphisms of $\bar k/k$, which proves Galois equivariance.

For precision, the prismatic input used for the other assertions is the following result of \citet[Theorems~3.3 and~4.14, Corollary~4.15, author-version pp.~15--16 and~27--28]{LiLiu}. For a smooth proper formal scheme over $\OO_K$, the $u$-power torsion module $M^j$ has finite length, and its Frobenius-fixed quotient after setting $u=0$ is exactly the kernel of specialization. The quotient is zero when $e(j-1)<p-1$. If $e(j-1)<2(p-1)$ it is killed by $p^{j-1}$. At equality $e(j-1)=p-1$, the module itself is killed by both $u$ and $p$ and has bijective Frobenius after the appropriate scalar extension. These assertions apply to the formal completion of the given smooth proper scheme; its geometric \'etale specializations are the groups appearing in $s^j$.

The role of the strict bound can be seen in the finite-length estimate behind that theorem. If the reduction of the annihilator of $M^j$ modulo $p$ is $(u^\alpha)$, its Frobenius-height relation gives
\[
 p\alpha\leq e(j-1)+\alpha.
\]
Under $e(j-1)<p-1$, this forces $\alpha=0$ and hence $M^j=0$. The next range permits only $\alpha\leq1$; the integral annihilator calculation in the cited theorem gives the stated $p^{j-1}$ bound and its sharper boundary form. The integral calculation, rather than a rational slope comparison, is needed here.

Taking Frobenius fixed points preserves the asserted annihilation bounds, and gives the three conclusions about $\ker s^j$. If the generic group is torsion-free and $s^j$ is injective, its special-fibre source is torsion-free as well. Equation~\eqref{eq:stable} now gives the last assertion. When $e=1$, the strict inequality is $j-1<p-1$, equivalently $j<p$.
\end{proof}

\begin{remark}\label{rem:scope}
The theorem concerns stable crystalline torsion, not all crystalline torsion. It gives a sufficient condition, not a converse. Outside the strict range, generic torsion-freeness merely forces any special \'etale torsion to lie in the specialization kernel; it does not force that kernel to vanish. The following construction realizes this possibility.
\end{remark}

\section{The integral K3 cover and the torsion section}

The following equations realize the cover and the disappearing torsion section needed in the quotient construction.

\begin{lemma}\label{lem:k3}
The scheme $\mathcal Y$ in \eqref{eq:quadrics} is a smooth projective K3 surface over $\ZZ_2$. The displayed involution is free, and its quotient $\mathcal S$ is smooth and projective. On every algebraically closed geometric fibre,
\[
 \pi_1^{\et}(\mathcal Y)=1,\qquad
 \pi_1^{\et}(\mathcal S)=\ZZ/2\ZZ.
\]
\end{lemma}
\begin{proof}
Substituting $x_i-y_i$ for $y_i$ leaves $y_i^2-x_iy_i$ unchanged, and applying the substitution twice is the identity. Thus $\tau$ is an integral involution preserving all three quadrics.

Work first over an algebraic closure of $\FF_2$. The Jacobian matrix has left block
\[
 \begin{pmatrix}
 y_1&x_3&x_2\\
 x_3&y_2&x_1\\
 x_2&x_1&y_3
 \end{pmatrix}
\]
and right block $\operatorname{diag}(x_1,x_2,x_3)$. If all $x_i$ are nonzero, the right block is invertible. Otherwise permute the indices so that $x_1=0$. If $x_2x_3\ne0$, the columns $dx_2,dy_2,dy_3$ have determinant $x_2x_3^2$, which is nonzero. The remaining case, after a further permutation, is $x_1=x_2=0$. The equations give $y_1=y_2=0$. The coordinate $x_3$ cannot vanish, since then $y_3=0$ and every coordinate would vanish. The columns $dx_1,dx_2,dy_3$ now have determinant $x_3^3\ne0$. Thus the Jacobian has rank three at every point.

Three hypersurfaces in $\PP^5$ have a nonempty intersection of dimension at least two. The rank calculation shows that this intersection is a smooth complete intersection of dimension exactly two. Its Koszul resolution by line bundles on $\PP^5$ gives $H^0(\OO)=k$ and $H^1(\OO)=0$: the required intermediate cohomology of $\OO_{\PP^5}(-2)$, $\OO_{\PP^5}(-4)$, and $\OO_{\PP^5}(-6)$ vanishes. Adjunction gives
\[
 \omega_{\mathcal Y_k}=\OO_{\mathcal Y_k}(-6+2+2+2)=\OO_{\mathcal Y_k}.
\]
In particular the special fibre is a geometrically connected K3 surface.

The relative Jacobian criterion makes the integral model smooth near its special fibre. The nonsmooth locus is closed in the projective $\ZZ_2$-scheme $\mathcal Y$, and its image in $\Spec\ZZ_2$ is closed. It misses the closed point, so it is empty. Thus $\mathcal Y$ is smooth over $\ZZ_2$. The same Koszul and adjunction calculations apply to its geometric generic fibre.

A projective fixed point of $\tau$ in characteristic two is represented by an eigenvector of the corresponding linear transformation. That transformation is unipotent and has only eigenvalue one. The fixed-vector equation gives every $x_i=0$, and the three quadrics then give every $y_i=0$, a contradiction. The projective fixed subscheme therefore has empty special fibre. Its properness makes it empty over $\ZZ_2$ as well.

A free action of the constant finite \'etale group $\ZZ/2\ZZ$ on this projective scheme has a scheme quotient, and the quotient map is a finite \'etale torsor \citep[Lemma~66.14.2, Tag~07S7]{Stacks}. One can construct the quotient using invariant affine opens, or the section ring of an invariant ample bundle. The tensor product of an ample bundle with its translate has a canonical linearization and descends along the torsor. Its descent is ample, since ampleness descends under a finite surjective morphism. Properness descends as well, so the quotient is projective. Smoothness descends along the finite \'etale cover.

We spell out the Lefschetz input for the fundamental group. A complete intersection of dimension at least two in projective space over an algebraically closed field is simply connected in the \'etale sense. Indeed, use its regular sequence of hypersurfaces as a flag. Every intermediate scheme is a local complete intersection, hence Cohen--Macaulay; at the steps in question its closed local rings have depth at least three. They also satisfy the purity hypothesis in \citet[Expos\'e~XII, Corollary~3.5, original p.~152; reedition p.~120]{SGA2}. That corollary gives an isomorphism of fundamental groups at each step down to dimension two. Projective space has trivial fundamental group: for $\PP^1$ this follows from Riemann--Hurwitz for an \'etale cover, and the surjectivity clause of the same Lefschetz theorem gives the higher-dimensional case. Applying this to the three quadrics proves $\pi_1^{\et}(\mathcal Y)=1$ on both geometric fibres. The exact sequence of the connected \'etale double cover $\mathcal Y\to\mathcal S$ then gives $\pi_1^{\et}(\mathcal S)=\ZZ/2\ZZ$.
\end{proof}

\begin{lemma}\label{lem:elliptic}
The curve $\mathcal E$ is a smooth projective elliptic curve over $\ZZ_2$ with ordinary special fibre. The point $P$ is a section of exact order two on the generic fibre and reduces to the identity on the special fibre.
\end{lemma}
\begin{proof}
For the displayed Weierstrass equation,
\[
 b_2=1,\quad b_4=0,\quad b_6=4,\quad b_8=1,
 \qquad \Delta=-b_2^2b_8-27b_6^2=-433.
\]
The discriminant is a $2$-adic unit. The standard point at infinity is the identity, so the smooth projective cubic is an elliptic curve over $\ZZ_2$.

The reduction of $f(U)=U^3+U^2+64$ has exactly one odd root, namely $U=1$. Its derivative $3U^2+2U$ is odd there. Hensel's lemma gives a unique odd $u\in\ZZ_2$ with $f(u)=0$. The three homogeneous coordinates $2u,-u,8$ are integral and generate the unit ideal because $u$ is odd. They define a section. The homogeneous cubic is
\[
 Y^2Z+XYZ=X^3+Z^3.
\]
At the displayed coordinates its left side is $-8u^2$, whereas its right side is $8u^3+512=-8u^2$. Hence the section lies on $\mathcal E$.

On the generic fibre its affine coordinates are $(u/4,-u/8)$. The inverse map is $(x,y)\mapsto(x,-y-x)$, which fixes this point. The point is not the point at infinity on that fibre, so it has exact order two. Its special reduction is $[0:1:0]$, the identity.

Finally the invariant differential reduces to $dx/x$ in characteristic two. It is nonzero and regular, and the Cartier operator fixes a logarithmic differential. Thus Cartier is bijective on the one-dimensional space of regular differentials, so the special elliptic curve is ordinary. Equivalently its Hasse coefficient $a_1$ is the unit one. Its maximal pro-$2$ \'etale fundamental group is consequently $\ZZ_2$.
\end{proof}

\section{The quotient threefold and its two cohomologies}

First-cohomology characters determine the second-cohomology torsion through a coefficient boundary.

\begin{lemma}\label{lem:bockstein}
For a connected proper scheme $V$ over an algebraically closed field and the usual integral $p$-adic \'etale cohomology, the coefficient sequence gives
\[
 0\longrightarrow H^1_{\et}(V,\ZZ_p)/p
 \longrightarrow H^1_{\et}(V,\FF_p)
 \longrightarrow H^2_{\et}(V,\ZZ_p)[p]
 \longrightarrow0.
\]
\end{lemma}
\begin{proof}
Take the long exact cohomology sequence of multiplication by $p$ on integral \'etale coefficients. Exactness at the three displayed terms gives the assertion. Equivalently take compatible finite coefficient sequences: their finite cohomology groups satisfy the Mittag--Leffler condition, so no additional inverse-limit term occurs in this degree. The character description of first cohomology is the classification of torsors under the constant finite cyclic groups, followed by the inverse limit.
\end{proof}

Let $\mathcal D$ be the quotient in Theorem~\ref{thm:surface}. The diagonal action is free because its action on $\mathcal Y$ is free. As in Lemma~\ref{lem:k3}, the quotient exists, a norm of an ample bundle proves projectivity, and finite \'etale descent proves smoothness. Its relative dimension is three. On the special fibre the elliptic translation is the identity, so
\begin{equation}\label{eq:product}
 D_0=S_0\times E_0.
\end{equation}

Over $\overline{\FF}_2$, the proper product formula for fundamental groups and the preceding lemmas give, on taking maximal pro-$2$ abelian quotients,
\[
 \pi_1^{\et}(D_0)^{\mathrm{ab},(2)}
       =(\ZZ/2\ZZ)\times\ZZ_2.
\]
The product formula here can also be obtained from the proper homotopy exact sequence applied to the projection in \eqref{eq:product}: the section splits the sequence, the other projection splits the inclusion of the fibre, and the product has trivial monodromy. Therefore
\begin{align*}
 H^1_{\et}(D_{0,\overline{\FF}_2},\ZZ_2)&=\ZZ_2,\\
 H^1_{\et}(D_{0,\overline{\FF}_2},\FF_2)&=\FF_2^2.
\end{align*}
The reduction of the first group is exactly the character space of the elliptic factor. Lemma~\ref{lem:bockstein} gives
\begin{equation}\label{eq:special-torsion}
 H^2_{\et}(D_{0,\overline{\FF}_2},\ZZ_2)[2]=\FF_2.
\end{equation}

On the geometric generic fibre $P$ is a nonzero two-torsion point. Let $E'=E/\langle P\rangle$, an elliptic curve over $\overline{\mathbb Q}_2$. The projection to $E$ followed by this isogeny descends to
\[
 D\longrightarrow E'.
\]
After the finite \'etale base change $E\to E'$, this morphism becomes the projection $Y\times E\to E$. Its geometric fibres are the simply connected K3 surface $Y$. It is smooth and proper with connected fibres. The proper homotopy exact sequence \citep[Proposition~58.15.2, Tag~0C0J]{Stacks} therefore gives
\[
 \pi_1^{\et}(D)\simeq\pi_1^{\et}(E').
\]
In characteristic zero the maximal pro-$2$ quotient of an elliptic fundamental group is $\ZZ_2^2$. Hence
\[
 H^1_{\et}(D,\ZZ_2)=\ZZ_2^2,
 \qquad H^1_{\et}(D,\FF_2)=\FF_2^2,
\]
and reduction is onto. Lemma~\ref{lem:bockstein} shows that $H^2_{\et}(D,\ZZ_2)[2]=0$. Since this is a finite $\ZZ_2$-module, it is torsion-free. This recovers the generic-fibre calculation in \citet[proof of Proposition~2.2, pp.~237--238]{BMS} while keeping the special integral obstruction visible.

\section{Cutting out a surface}

An ample hypersurface preserves the fundamental groups used above and therefore preserves the decisive coefficient boundary.

\begin{proof}[Proof of Theorem~\ref{thm:surface}]
The construction and smooth projectivity of $\mathcal D$ were established above. Choose a projective embedding over $\ZZ_2$ given by a sufficiently high power of an ample bundle. Its special fibre is smooth and geometrically connected of dimension three. Poonen's finite-field Bertini theorem says that smooth hypersurface sections of dimension two have positive asymptotic density among homogeneous equations of large degree \citep[Theorem~1.1 and its proof in \S2.4, author-version pp.~2 and~7]{Poonen}. Choose one such equation over $\FF_2$.

Lift its coefficients to $0$ or $1$ in $\ZZ_2$, and let $\mathcal X$ be its intersection with $\mathcal D$. The special intersection is a smooth effective ample Cartier divisor $X_0$. The relative Jacobian criterion shows that $\mathcal X$ is smooth near $X_0$. As before, the nonsmooth locus is proper over $\ZZ_2$ and has empty special fibre, so it is empty. The resulting $\mathcal X$ is a smooth projective relative surface.

Over each algebraically closed geometric field, the connectedness and fundamental-group clauses of \citet[Expos\'e~XII, Corollary~3.5]{SGA2} apply to the ample divisor in the smooth threefold. Indeed the ambient local rings have depth three at closed points and are regular away from the divisor, so its depth and purity hypotheses are satisfied. Thus both geometric fibres of $\mathcal X$ are connected and
\[
 \pi_1^{\et}(X_{0,\overline{\FF}_2})
    \simeq\pi_1^{\et}(D_{0,\overline{\FF}_2}),
 \qquad
 \pi_1^{\et}(\mathcal X_{\overline{\mathbb Q}_2})
    \simeq\pi_1^{\et}(\mathcal D_{\overline{\mathbb Q}_2}).
\]
These identifications respect both integral characters and their reduction modulo two. Lemma~\ref{lem:bockstein} and the preceding section therefore imply
\[
 H^2_{\et}(X_{0,\overline{\FF}_2},\ZZ_2)[2]=\FF_2,
 \qquad
 H^2_{\et}(\mathcal X_{\overline{\mathbb Q}_2},\ZZ_2)[2]=0.
\]
The generic group is finite over $\ZZ_2$, so the second equality makes it torsion-free. Equation~\eqref{eq:stable}, applied to the smooth projective special surface, turns the first equality into
\[
 T^2(X_{0,\overline{\FF}_2})[2]_{\st}
      =\FF_2\otimes_{\FF_2}\overline{\FF}_2
      =\overline{\FF}_2.
\]
The absolute ramification index is one and $j=p=2$. Thus $e(j-1)=p-1$, exactly the first excluded equality in the strict vanishing range.
\end{proof}

\begin{remark}\label{rem:effective}
Only the last hypersurface is selected by existence. It can be chosen by enumerating homogeneous equations over $\FF_2$ by degree and testing geometric smoothness of their intersections using the Jacobian ideal. Poonen's theorem proves termination. This is not an explicit degree bound. The equations for the K3 cover, its involution, the elliptic curve, and its torsion section are explicit, and the Jacobian and discriminant calculations above are exact. The program \emph{Integral lifting arithmetic}, available through \url{\workbenchurl}, checks these polynomial identities. No finite numerical search is used to justify the cohomological conclusions.
\end{remark}

\section{What the example does and does not say}

The failure is caused by a special-fibre \'etale character which has no integral generic counterpart. In \eqref{eq:product} it comes from the connected double cover of the Enriques quotient. Generically that cover is coupled to a nonzero elliptic two-torsion translation, and the resulting threefold has the same fundamental group as an elliptic curve. The coefficient exact sequence measures the lost character as second-cohomology torsion. Ogus then identifies precisely its Frobenius-stable crystalline part.

This explanation is compatible with the comparison results of Bhatt--Morrow--Scholze: generic \'etale torsion-freeness does not imply crystalline torsion-freeness in unrestricted degree and ramification. Their construction already supplies that counterexample. Li--Liu describe the specialization kernel which accounts for it and show why the strict range behaves differently. The explicit model here therefore resolves the proposed unrestricted implication negatively without asserting a new counterexample phenomenon or a new comparison theorem.

The construction does not classify specialization kernels beyond the stated range. It does not show that all crystalline torsion is stable, nor that stable torsion is determined by the generic rational cohomology. Those are different assertions, and no such assertion enters either proof.

\section*{Note on method and reuse}

The mathematics and text were produced by GPT-6 Astra in ChatGPT Pro mode, whose displayed model is GPT-6 Astra Pro. The explicit geometry and cohomological reductions were checked against the cited primary passages and by separate exact arithmetic calculations. The known comparison theorems remain attributed to their authors. A separate instance of the same model reconstructed both main assertions from their statements before comparison with this text. That reconstruction used the explicitly cited fundamental comparison theorems as published inputs and gave a separate local Artin--Schreier verification of the integral fixed-point identity. The complete attempts were frozen before comparison, and their geometric and cohomological deductions were checked in full. Reasoning-effort metadata is not exposed. No human has checked the mathematics or text. The text is dedicated to the public domain under CC0~1.0. The proofs, check records, and editable source belong to the workbench at \url{\workbenchurl}; the permanent record address is \url{\recordurl}.

\begin{thebibliography}{9}
\bibitem[Bhatt et~al.(2018)]{BMS}
Bhatt, B., Morrow, M., \& Scholze, P. (2018). Integral $p$-adic Hodge theory. \emph{Publications Math\'ematiques de l'IH\'ES, 128}, 219--397. \url{https://doi.org/10.1007/s10240-019-00102-z}.
\bibitem[Grothendieck(n.d.)]{Grothendieck}
Grothendieck, A. (n.d.). \emph{Torsion en cohomologie cristalline, cote~4} [Unpublished manuscript]. Archives Grothendieck, Universit\'e de Montpellier. \S8.5.10, archive sheet~9, scan image~10. \url{https://grothendieck.umontpellier.fr/4.pdf}.
\bibitem[Grothendieck(2005)]{SGA2}
Grothendieck, A. (2005). \emph{Cohomologie locale des faisceaux coh\'erents et th\'eor\`emes de Lefschetz locaux et globaux (SGA~2)} (Y. Laszlo, Ed.; reedition, Documents Math\'ematiques~4). Soci\'et\'e Math\'ematique de France. Original seminar, 1962; first edition, 1968. Expos\'e~XII, Corollary~3.5. \url{https://www.cmls.polytechnique.fr/perso/laszlo/sga2/sga2-smf.pdf}.
\bibitem[Li and Liu(2023)]{LiLiu}
Li, S., \& Liu, T. (2023). On the $u^\infty$-torsion submodule of prismatic cohomology. \emph{Compositio Mathematica, 159}(8), 1607--1672. \url{https://doi.org/10.1112/S0010437X23007261}. Exact theorem and page locators used here refer to the verified author version: \url{https://arxiv.org/abs/2204.03411v1}.
\bibitem[Ogus(1975)]{Ogus}
Ogus, A. (1975). Cohomology of the infinitesimal site. \emph{Annales scientifiques de l'\'Ecole normale sup\'erieure, 8}(3), 295--318. \url{https://doi.org/10.24033/asens.1289}.
\bibitem[OpenAI(2026)]{OpenAI}
OpenAI. (2026, October~7). \emph{Milne's rationality conjecture for abelian varieties} [Preprint]. \url{https://github.com/openai/math/tree/main/preprints/Milnes-rationality-conjecture-for-abelian-varieties-October-7-2026}.
\bibitem[Poonen(2004)]{Poonen}
Poonen, B. (2004). Bertini theorems over finite fields. \emph{Annals of Mathematics, 160}(3), 1099--1127. \url{https://doi.org/10.4007/annals.2004.160.1099}. Author version: \url{https://math.mit.edu/~poonen/papers/bertini.pdf}.
\bibitem[The Stacks Project Authors(2026)]{Stacks}
The Stacks Project Authors. (2026). \emph{The Stacks Project}. Tags~07S7 and~0C0J. \url{https://stacks.math.columbia.edu}.
\bibitem[Zou(2026)]{Zou}
Zou, H. (2026, November~3). \emph{Bogomolov--Tian--Todorov theorems via crystalline periods of differential forms} [Draft preprint]. \url{https://www.math.uni-bielefeld.de/~hzou/papers/btt/}.
\end{thebibliography}
\end{document}
